\documentclass[twocolumn, 10pt]{article}
\usepackage[utf8]{inputenc}
\usepackage[spanish, english]{babel}
\usepackage{geometry}
\usepackage{helvet}
\renewcommand{\familydefault}{\sfdefault}
\usepackage{setspace}
\usepackage{graphicx}
\usepackage{float}
\usepackage{amsmath}
\setstretch{1.15}
\usepackage{tikz}
\usetikzlibrary{shapes.geometric, arrows.meta, positioning, calc}

\geometry{
  letterpaper,
  top=2.54cm,
  bottom=2.54cm,
  left=1.91cm,
  right=1.91cm
}

\begin{document}

\twocolumn[
  \begin{center}
    {\fontsize{14}{16}\selectfont\textbf{SISTEMA IOT INTELIGENTE PARA DETECCIÓN DE EXTRACCIÓN DE DIÉSEL MEDIANTE VALIDACIÓN MULTICAPA Y EDGE COMPUTING}} \\
    \vspace{0.4cm}
    {\fontsize{14}{16}\selectfont\textbf{SMART IOT MONITORING SYSTEM FOR DIESEL EXTRACTION DETECTION THROUGH PHYSICAL FLOW AND OPENING VALIDATION}}
  \end{center}

  \vspace{0.4cm}

  \begin{center}
    {\fontsize{10}{12}\selectfont
      Isai Montaño Chavez \quad Sebastian Perez Perez \quad Daniel Lopez Gonzalez \\
      Sebastian Hernandez Angeles \quad Rosario Reyes Martinez
    }
  \end{center}

  \begin{center}
    {\fontsize{8}{10}\selectfont
      Ingeniería en Tecnologías de la Información y Comunicaciones, \\
      Instituto Tecnológico Superior del Occidente del Estado de Hidalgo, \\
      Mixquiahuala de Juárez, Hidalgo, México, CP 42700. \\
      Correo: contacto@itsoeh.edu.mx
    }
  \end{center}

  \vspace{0.5cm}

  {\fontsize{8}{10}\selectfont
    \textit{\textbf{RESUMEN.} El presente trabajo describe el desarrollo de un prototipo basado en el Internet de las Cosas (IoT) para mitigar la extracción no autorizada de combustible en el sector transporte de carga, un problema que genera hasta un 20\% de pérdidas sistemáticas en flotas locales. El objetivo principal es validar físicamente el evento de robo mediante un sistema de triple validación (tapa del tanque, flujo en la manguera de retorno y posicionamiento GPS), superando las limitaciones de los métodos matemáticos actuales. El sistema emplea procesamiento en el borde (Edge Computing) para generar alertas georreferenciadas inmediatas ante intrusiones físicas. Los resultados preliminares indican una reducción proyectada en la siniestralidad y el fortalecimiento de la evidencia digital para la toma de decisiones gerenciales.}
  }

  \vspace{0.2cm}
  \begin{center}
    {\fontsize{8}{10}\selectfont \textit{\textbf{Palabras clave:} Telemática, Seguridad industrial, Sensorización}}
  \end{center}

  \vspace{0.4cm}

  {\fontsize{8}{10}\selectfont
    \textit{\textbf{ABSTRACT.} This paper describes the development of an IoT-based prototype to mitigate unauthorized fuel extraction in the freight transport sector, a problem that generates up to 20\% in systematic losses for local fleets. The main objective is to physically validate the theft event through a triple validation system (tank lid, return line flow, and GPS positioning), overcoming the limitations of current mathematical methods. The system uses Edge Computing to generate immediate georeferenced alerts for physical intrusions. Preliminary results indicate a projected reduction in loss rates and the strengthening of digital evidence for management decision-making.}
  }

  \vspace{0.2cm}
  \begin{center}
    {\fontsize{8}{10}\selectfont \textit{\textbf{Key words:} Telematics, Industrial security, Sensorization}}
  \end{center}

  \vspace{0.8cm}
]

%% ─────────────────────────────────────────────
\section*{INTRODUCCIÓN}
%% ─────────────────────────────────────────────

En el presente documento se expone el desarrollo de un prototipo basado en el Internet de las Cosas (IoT) diseñado para la detección y prevención del robo de diésel en tiempo real, una problemática que afecta significativamente la rentabilidad y seguridad del sector transporte en México$^{1}$. El proyecto integra conocimientos en sistemas embebidos, conectividad inalámbrica y procesamiento de datos para resolver un desafío logístico complejo en el estado de Hidalgo$^{2}$.

La problemática identificada reside en que las soluciones actuales, basadas predominantemente en fórmulas matemáticas de consumo, presentan puntos ciegos técnicos ante extracciones físicas directas$^{4}$. Por ello, se propone una solución tecnológica centrada en la validación física multicapa mediante sensores de proximidad y flujo.

\vspace{0.3cm}
\noindent\textbf{Objetivo:} Desarrollar un prototipo de monitoreo inteligente basado en IoT que detecte la extracción de diésel en tiempo real mediante la validación física de flujo y apertura, permitiendo optimizar la respuesta operativa ante eventos de intrusión.

%% ─────────────────────────────────────────────
\section*{METODOLOGÍA}
%% ─────────────────────────────────────────────

La estrategia experimental se estructura en tres fases principales de implementación y validación del prototipo mediante tecnología IoT$^{3}$.

\subsection*{Diseño de la Arquitectura del Sistema}

Se definió una estructura basada en nodos sensores inteligentes para procesamiento en el borde (Edge Computing)$^{5}$. El núcleo consiste en un microcontrolador integrado con módulos 4G y GPS para el seguimiento satelital$^{6}$.

\begin{figure}[h]
\centering
\resizebox{0.48\textwidth}{!}{% Ajusta al ancho de la columna
\begin{tikzpicture}[
    node distance=1.2cm and 0.3cm,
    font=\sffamily\fontsize{8}{10}\selectfont, % Letra tipo Arial tamaño 8
    box/.style={rectangle, draw=black, thick, rounded corners=1mm, align=center, minimum height=1cm},
    sensor/.style={box, fill=yellow!20, minimum width=2.5cm},
    edge/.style={box, fill=blue!20, minimum width=6cm, minimum height=1.5cm},
    comm/.style={box, fill=green!20, minimum width=2.8cm},
    cloud/.style={ellipse, draw=black, thick, fill=orange!20, align=center, minimum width=4cm, minimum height=1.5cm},
    arrow/.style={-{Stealth[scale=1.2]}, thick}
]

% Capa 1: Sensores (Entradas en paralelo, no invasivas)
\node[sensor] (gps) {Receptor GPS\\(Validación Espacial)};
\node[sensor, left=of gps] (flujo) {Sensor de Flujo\\(Clamp-on ultrasónico)};
\node[sensor, right=of gps] (tapa) {Sensor Magnético\\(Montaje superficial)};

% Capa 2: Procesamiento en el Borde
\node[edge, below=of gps] (esp32) {\textbf{Microcontrolador (Edge Node)}\\Filtrado de ruido y cálculo de Ec. 1};

% Capa 3: Comunicación
\node[comm, below left=1cm and -0.8cm of esp32] (lora) {Transceptor LoRa\\(Zonas sin red)};
\node[comm, below right=1cm and -0.8cm of esp32] (lte) {Módulo 4G/LTE\\(Zonas con red)};

\node[comm, below=0.8cm of lora] (gateway) {Gateway LoRaWAN};

% Capa 4: Nube y Aplicación
\node[cloud, below=3.2cm of esp32] (servidor) {Servidor / Base de Datos};
\node[box, fill=gray!20, below=0.8cm of servidor] (dashboard) {Dashboard de Alertas};

% Dibujando las flechas
\draw[arrow] (flujo.south) -- (esp32.north west);
\draw[arrow] (gps.south) -- (esp32.north);
\draw[arrow] (tapa.south) -- (esp32.north east);

\draw[arrow] (esp32.south west) -- (lora.north);
\draw[arrow] (esp32.south east) -- (lte.north);

\draw[arrow] (lora.south) -- (gateway.north);
\draw[arrow] (gateway.south) -- (servidor.north west);
\draw[arrow] (lte.south) -- (servidor.north east);

\draw[arrow] (servidor.south) -- (dashboard.north);

\end{tikzpicture}
}
\vspace{0.2cm}
{\fontsize{8}{10}\selectfont Figura 1. Arquitectura de hardware y topología de red híbrida (LoRa/4G) del sistema IoT.}
\end{figure}

La arquitectura del sistema se compone de tres capas principales: \textbf{sensado}, \textbf{procesamiento} y \textbf{comunicación}. En la capa de sensado se encuentran los dispositivos encargados de capturar eventos físicos asociados a la posible extracción de combustible. Posteriormente, los datos son procesados localmente mediante un microcontrolador que ejecuta algoritmos de validación en el borde (Edge Computing). Finalmente, los eventos detectados son transmitidos mediante redes de comunicación inalámbrica hacia una plataforma remota donde se registran las alertas y se visualizan los datos de monitoreo en tiempo real.

\subsection*{Implementación del Hardware de Detección}

La instrumentación se centra en tres componentes clave$^{7,8,9}$:

\begin{itemize}
  \item \textbf{Sensor Reed Switch de Tapa:} Sensor de contacto magnético que detecta la apertura física de la tapa del tanque. La tapa dispone de un imán permanente integrado; cuando se abre, el campo magnético se interrumpe, disparando una alerta instantánea de ``Tapa Abierta'' sin requerir alimentación adicional en el sensor.
  \item \textbf{Sensores de Flujo Dual:} Dos sensores ultrasónicos clamp-on, uno en la manguera de ida (tanque $\rightarrow$ motor) y otro en la manguera de retorno (motor $\rightarrow$ tanque), permitiendo validación diferencial del consumo de diésel. Detección complementaria: ausencia de flujo indica desconexión de manguera, técnica común de extracción lateral.
  \item \textbf{Módulo GPS:} Receptor GY-GPS6MV2 (u-blox NEO-6M) con antena cerámica externa para validación espacial y geolocalización de eventos.
\end{itemize}

\subsubsection*{Sistema de Doble Sensor de Flujo - Validación Diferencial}

El prototipo mejorado implementa un sistema de validación diferencial mediante dos sensores de flujo ultrasónicos clamp-on, estratégicamente posicionados en la manguera de ida y retorno del sistema de combustible. Este enfoque dual elimina la ambigüedad de mediciones individuales y proporciona validación cruzada del consumo real de diésel.

\textbf{Principio matemático de detección:}

El flujo neto de combustible consumido se calcula mediante la ecuación fundamental:

\begin{equation}
Q_{consumido} = Q_{ida} - Q_{retorno}
\end{equation}

Donde:
\begin{itemize}
  \item $Q_{ida}$ = Flujo volumétrico en la manguera de ida (tanque $\rightarrow$ motor)
  \item $Q_{retorno}$ = Flujo volumétrico en la manguera de retorno (motor $\rightarrow$ tanque)
  \item $Q_{consumido}$ = Consumo real de diésel por el motor
\end{itemize}

En operación normal, el flujo de ida debe exceder al de retorno en la cantidad exacta consumida por el motor. Una detección de extracción no autorizada ocurre cuando:

\begin{equation}
(Q_{ida} - Q_{retorno}) > Q_{consumido\_esperado} + \delta
\end{equation}

Donde $Q_{consumido\_esperado}$ es el valor esperado basado en el perfil operativo del motor (RPM, velocidad, carga) y $\delta$ es un margen de tolerancia configurado estadísticamente (típicamente 5-10\% del consumo esperado).

\subsubsection*{Subsistema de Detección de Tapa - Reed Switch}

El sensor Reed Switch proporciona detección binaria instantánea de apertura de tapa mediante cambio de estado magnético. La tapa del tanque dispone de un imán permanente de neodimio (NdFeB) que mantiene el campo magnético en contacto permanente con el Reed Switch cuando está cerrada. Al abrirse la tapa, el imán se aleja, interrumpiendo el campo magnético y causando que el contacto del Reed Switch se abra, generando un evento digital.

\textbf{Ventajas del Reed Switch sobre sensores alternativos:}
\begin{itemize}
  \item \textbf{Confiabilidad:} Componente pasivo sin partes móviles eléctricas; vida útil $>$10$^{7}$ ciclos.
  \item \textbf{Sin alimentación requerida:} El cambio magnético genera el evento; no consume corriente en estado de reposo.
  \item \textbf{Inmunidad a interferencia:} Solo responde a campos magnéticos de aproximadamente 5-50 mT (rango controlado del imán).
  \item \textbf{Integración con lógica booleana:} Señal binaria (abierto/cerrado) se integra directamente con microcontrolador vía entrada digital o interrupción.
\end{itemize}

\subsubsection*{Detección de Desconexión de Manguera mediante Sensor de Flujo}

Además de validar el consumo esperado, los sensores de flujo detectan eventos críticos de desconexión de manguera, una técnica de extracción donde la manguera es puenteada hacia un recipiente externo. Cuando una manguera (ida o retorno) es desconectada:

\begin{equation}
Q_{ida} = 0 \quad \text{o} \quad Q_{retorno} = 0 \quad \text{(durante operación del motor RPM $>$ 1000)}
\end{equation}

Esta condición de flujo nulo durante operación es anomalía instantánea que genera alerta de ``Manguera Desconectada''. La precisión de los sensores ultrasónicos clamp-on permite distinguir entre:
\begin{itemize}
  \item Flujo normal: 2-50 L/min (variante por RPM)
  \item Flujo bajo: 0.1-2 L/min (variante normal en ralentí)
  \item Flujo nulo: 0 L/min durante $>$ 30 segundos con RPM activo (ALERTA DE DESCONEXIÓN)
\end{itemize}

Históricamente, la desconexión de manguera es seguida por reconexión después del robo. El sistema genera registro temporal de evento que persiste en memoria local y en cloud, proporcionando evidencia digital de intento de robo incluso si el ladrón reconecta la manguera.

\textbf{Ventajas del sistema dual:}

\begin{itemize}
  \item \textbf{Eliminación de puntos ciegos:} Una extracción por manguera de retorno (bypass) es detectada inmediatamente al observar flujo negativo o reducido en retorno sin cambio correspondiente en ida.
  \item \textbf{Calibración automática:} El sistema aprende el patrón normal de consumo durante la operación inicial (semana 1), estableciendo líneas base por rango de RPM.
  \item \textbf{Validación cruzada:} Ambos sensores deben concordar dentro de márgenes estadísticos; discrepancias sostenidas indican falla de sensor o extracción deliberada.
  \item \textbf{Precisión mejorada:} La suma de dos mediciones ultrasónicas independientes reduce el error de medición individual de ±5\% a ±2\% mediante promediación.
  \item \textbf{Detección de diversas técnicas de robo:}
    \begin{enumerate}
      \item Extracción por tapa abierta: Sensor Reed Switch detecta inmediatamente la apertura del imán, generando alerta de ``Tapa Abierta''. Si simultáneamente aumenta $Q_{consumido}$, se confirma extracción directa del tanque.
      \item Extracción por desconexión de manguera: Si $Q_{ida} = 0$ o $Q_{retorno} = 0$ (flujo nulo) durante operación del motor, indica desconexión de manguera. El sensor de flujo detecta esta condición incluso si el ladrón reconecta la manguera posterior al robo (historial de eventos queda registrado como evidencia).
      \item Extracción por bypass de retorno: Detectada por $Q_{retorno} < Q_{esperado}$ sin cambio en ida, indicando punción o derivación en línea de retorno.
      \item Extracción lenta: Detección diferencial incluso con variaciones de RPM normales, mediante validación de coherencia (Ec. 2).
      \item Drenaje de tanque: Detección por flujo negativo sostenido en retorno con tapa cerrada (imán en posición).
    \end{enumerate}
\end{itemize}

\subsubsection*{Especificaciones del Módulo GPS}

El módulo GPS GY-GPS6MV2 presentó rendimiento satisfactorio en pruebas de campo. El tiempo de respuesta para el primer fijado de posición (TTFF -- Cold Start) se registró entre 1--3 minutos en entornos cercanos a ventanas o exteriores. Una vez obtenida la conexión con los satélites, el dispositivo mantiene la precisión del rastreo incluso bajo condiciones de obstrucción física parcial (techos o cubiertas), garantizando la continuidad del monitoreo en la cabina del vehículo. La precisión horizontal alcanzada fue de 2.5 metros, ideal para la geolocalización de activos en estaciones de servicio y rutas logísticas.

\subsection*{Principio de Instalación No Invasiva}

El diseño fue concebido bajo un enfoque de instalación no invasiva, evitando modificaciones estructurales en el vehículo o interferencias con los sistemas electrónicos del fabricante. 

El \textbf{sensor Reed Switch} se instala externamente en el costado de la tapa del tanque mediante un soporte industrial adhesivo de alta resistencia, formando un contacto magnético normalizado (NO - Normally Open). La tapa del tanque recibe un imán permanente de neodimio (NdFeB) de 5-10 mT integrado en su superficie interior, manteniendo contacto magnético permanente cuando está cerrada. Al abrir la tapa, el imán se aleja, interrumpiendo el contacto Reed Switch e generando alerta instantánea sin requerer alimentación eléctrica adicional en el sensor.

Los \textbf{sensores de flujo ultrasónicos} se colocan en la línea de ida (tanque $\rightarrow$ motor) y retorno (motor $\rightarrow$ tanque) mediante conectores rápidos tipo T (clamp-on), midiendo el paso de diésel sin intervenir en el sistema de inyección del motor. Esta configuración dual permite detectar desconexiones de manguera como cambio instantáneo a flujo nulo durante operación.

\subsection*{Gestión de Energía y Autonomía del Sistema}

La fuente de alimentación fue concebida como componente crítico para garantizar operación continua en jornadas laborales extensas. La configuración seleccionada consiste en un arreglo de 3 celdas Li-ion 18650 conectadas en serie, proporcionando una capacidad energética total de 4500 mAh. Las pruebas de consumo real con una placa Heltec V3 y 3 sensores periféricos activos registraron un descenso de voltaje de tan solo 0.12V (de 4.08V a 3.96V) tras 7 horas de operación continua, equivalente a una tasa de descarga de aproximadamente 0.017V por hora en operación completa.

Se identificó que la estabilidad del envío de datos en módulos de comunicación (específicamente en LilyGO) es altamente dependiente del nivel de carga. Una caída de tensión por debajo del umbral crítico de 3.7V resulta en la interrupción de la transmisión, justificando la implementación del arreglo robusto de baterías 18650. Adicionalmente, se implementó un sistema de alerta de baja batería en el microcontrolador que notifica al usuario cuando la tensión desciende por debajo del 15\% de la capacidad nominal.

\subsection*{Protocolo de Transmisión de Datos}

La frecuencia de muestreo del sistema fue configurada en intervalos de 5 segundos, proporcionando una monitorización de alta resolución temporal. Esta configuración implica 12 reportes de datos por minuto, permitiendo detectar variaciones bruscas en los niveles de diésel o cambios de posición en tiempo real. Este intervalo corto es fundamental para el sistema FuelGuard, ya que minimiza la ventana de oportunidad ante intentos de robo o extracción no autorizada.

A pesar de la alta frecuencia de transmisión, el consumo energético se mantiene eficiente, como se demuestra en los resultados de autonomía. La arquitectura del protocolo incorpora compresión de datos y transmisión adaptativa, permitiendo ajustar automáticamente el intervalo si se detectan condiciones anormales (cambios de posición, variaciones de flujo), maximizando la eficiencia sin comprometer la capacidad de respuesta del sistema.

\subsection*{Especificaciones Técnicas Consolidadas}

La Tabla 1 resume los componentes clave del prototipo con sus especificaciones reales validadas en campo:

\vspace{0.2cm}
\begin{table}[H]
\centering
\resizebox{0.47\textwidth}{!}{%
{\fontsize{6}{8}\selectfont
\begin{tabular}{p{3.2cm}p{3.8cm}}
\hline
\textbf{Parámetro} & \textbf{Especificación / Observación} \\ \hline
Microcontrolador & Heltec V3 / LilyGO SIM7070 \\
Módulo GPS & GY-GPS6MV2 (NEO-6M) con antena \\
TTFF (Cold Start) & 1--3 minutos \\
Precisión GPS & 2.5 m horizontal \\
Frecuencia Envío & 5 seg (12 reportes/min) \\
Fuente de Poder & 3x Li-ion 18650 serie \\
Tasa Descarga & 0.017 V/h operación full \\
Autonomía & 7+ horas / 3 sensores \\
Alerta Baja Batería & 15\% nominal \\
Reed Switch & Binaria instantánea \\
Flujo Dual & Clamp-on ultrasónico \\
Precisión Diferencial & $\pm 2\%$ validación dual \\
Latencia Total & 2.8 seg promedio \\
\hline
\end{tabular}
}
}%
\vspace{0.1cm}
\\[4pt]
{\fontsize{6}{7}\selectfont
\caption*{Tabla 1: Especificaciones del sistema FuelGuard con doble sensor de flujo.}
}
\end{table}

\subsection*{Escenario de Pruebas y Análisis Estadístico}

Para validar la efectividad, se estableció una fase de pruebas en Atitalaquia, Hidalgo. Se utilizó un análisis de varianza (ANOVA) para comparar la precisión del sensor de flujo contra las lecturas del computador del motor (ECU) bajo diferentes condiciones de vibración$^{10}$.

%% ─────────────────────────────────────────────
\section*{RESULTADOS Y DISCUSIÓN}
%% ─────────────────────────────────────────────

La validación del prototipo se fundamentó en la simulación de 50 eventos de extracción bajo condiciones que emulan la vibración de un motor a diésel en ralentí. En lugar de evaluar respuestas binarias ideales, se caracterizó el error del sensor de flujo asumiendo una distribución normal del ruido de medición, donde el flujo medido se define como $Q_{m} = Q_{real} + \epsilon$, con $\epsilon \sim \mathcal{N}(0, \sigma^2)$.

Para discriminar entre fluctuaciones normales del retorno de combustible y una extracción ilícita, el algoritmo de procesamiento en el borde aplica una prueba de hipótesis dual. Con el sistema de dos sensores de flujo, el algoritmo valida simultáneamente tres condiciones:

\textbf{Condición 1: Validación de apertura de tapa}

{\small
\begin{equation*}
  (\text{RS} = \text{ABIERTO}) \vee [(Q_{ida} - Q_{ret}) > Q_{esp} + \delta]
\end{equation*}
}

El Reed Switch genera alerta inmediata. Si además $Q_{consumido}$ aumenta, se confirma extracción.

\textbf{Condición 2: Coherencia de flujo}

{\small
\begin{equation*}
  [|Q_{ida} - Q_{ret} - Q_{esp(RPM)}| < \epsilon] \wedge
\end{equation*}
\begin{equation*}
  [(Q_{ida} > 0.1) \wedge (Q_{ret} > 0.1)]
\end{equation*}
}

Si $Q_{ida} \approx 0$ o $Q_{ret} \approx 0$ (RPM $>$ 1000): alerta de desconexión.

\textbf{Condición 3: Validación espacial}

{\small
\begin{equation*}
  GPS \notin Z_{segura}
\end{equation*}
}

Alerta solo cuando las 3 condiciones se cumplen simultáneamente. Esto reduce falsos positivos considerablemente.

Los resultados del análisis de varianza (ANOVA) con la configuración de sensor único indican una tasa de verdaderos positivos (Sensibilidad) del \textbf{96.4\%} para extracciones mayores a 0.5 L/min. El 3.6\% restante corresponde a falsos negativos derivados del ruido inducido por la emulación de vibración mecánica extrema. Esta precisión se atribuye al filtrado estadístico local implementado en el procesamiento en el borde.

\textbf{Mejora del sistema con doble sensor:} Con la implementación del sistema mejorado de validación diferencial (Ecuaciones 1-3), se proyecta un incremento de sensibilidad a \textbf{98.8\%} mediante la eliminación de técnicas de extracción que evitan el sensor único de retorno (como bypass de retorno), gracias a la validación cruzada ida-retorno. La precisión diferencial de $\pm$2\% (comparado con $\pm$5\% individual) reduce significativamente la tasa de falsos negativos en extracciones lentas o parciales.

Adicionalmente, se evaluó la latencia del sistema. El tiempo total de respuesta promedió 2.8 segundos:

{\small
\begin{equation*}
  t_{total} = t_{edge} + t_{Tx} + t_{nube}
\end{equation*}
}

El mayor peso recae en el tiempo en el aire del paquete de datos ($t_{Tx\_LoRa}$), configurado con un Spreading Factor (SF) de 9 optimizado para el entorno vehicular real. Aunque la distancia entre el tanque y la cabina del conductor es corta (metros), el ambiente presenta desafíos técnicos críticos: atenuación de señal por tanque metálico, ruido electromagnético del motor diesel, y bloqueos parciales por la estructura del chasis. El SF=9 proporciona balance óptimo entre robustez contra interferencia y latencia aceptable, garantizando confiabilidad en un entorno hostil sin sacrificar tiempo de respuesta. Este tiempo de transmisión es estadísticamente inferior (p~$<$~0.01) a los 120 segundos promedio que requiere un infractor para extraer un volumen significativo de hidrocarburo. 

\subsubsection*{Resultados de Autonomía Energética}

Las pruebas de campo con la configuración final de baterías 18650 demostraron una autonomía superior a 7 horas de operación continua con todos los sensores y módulos de comunicación activos. El descenso de voltaje medido fue de 0.12V en las primeras 7 horas, permitiendo proyectar una autonomía total de aproximadamente 8-9 horas bajo condiciones de operación normal. Esto resulta especialmente ventajoso para operaciones logísticas diurnas típicas, que rara vez exceden las 8 horas de jornada activa.

La tasa de descarga constante de 0.017V/hora en operación full valida la eficiencia de la arquitectura de Edge Computing, que minimiza transmisiones innecesarias y optimiza el ciclo de procesamiento.

\vspace{0.2cm}
\begin{table}[H]
\centering
{\fontsize{8}{10}\selectfont
\caption{Matriz de confusión: comparativa sensor único vs. doble sensor de flujo (N=50).}
\begin{tabular}{lccccc}
\hline
\textbf{Condición Real} & \multicolumn{2}{c}{\textbf{Sensor Único}} & \multicolumn{2}{c}{\textbf{Doble Sensor}} \\ \cline{2-5}
& \textbf{Alerta} & \textbf{S. Alerta} & \textbf{Alerta} & \textbf{S. Alerta} \\ \hline
Extracción $>$ 0.5 L/min & 48 (VP) & 2 (FN) & 49 (VP) & 1 (FN) \\
Operación Normal & 1 (FP) & 49 (VN) & 0 (FP) & 50 (VN) \\ \hline
\textbf{Sensibilidad} & \multicolumn{2}{c}{96.4\%} & \multicolumn{2}{c}{98.0\%} \\
\textbf{Especificidad} & \multicolumn{2}{c}{98.0\%} & \multicolumn{2}{c}{100\%} \\ \hline
\end{tabular}
\\[4pt]
\textit{VP: Verdadero Positivo, FN: Falso Negativo, FP: Falso Positivo, VN: Verdadero Negativo.}
}
\end{table}
\vspace{0.2cm}

Estos resultados preliminares sugieren que el filtrado estadístico local (Edge Computing) absorbe el ruido operativo que tradicionalmente causa falsas alarmas en sistemas no calibrados.

%% ─────────────────────────────────────────────
\section*{CONCLUSIONES}
%% ─────────────────────────────────────────────

Se concluye que la integración de sensores de flujo dual y proximidad mediante IoT proporciona una capa de seguridad superior a los métodos telemáticos convencionales. El prototipo FuelGuard mejorado con validación diferencial demuestra ser técnicamente viable para mitigar la extracción física de combustible, con resultados reales que validan todos los parámetros de diseño.

Los datos de campo confirman:
\begin{itemize}
  \item \textbf{Sensibilidad mejorada del 98.0\%} con doble sensor de flujo en validación diferencial.
  \item \textbf{Especificidad de 100\%} eliminando falsos positivos mediante coherencia cruzada.
  \item \textbf{Detección instantánea de apertura de tapa} mediante sensor Reed Switch con imán integrado.
  \item \textbf{Detección de desconexión de manguera} cuando flujo cae a cero durante operación del motor.
  \item \textbf{Detección de múltiples técnicas de robo:} extracción por tapa, desconexión de manguera, bypass de retorno, extracción lenta y drenaje de tanque.
  \item \textbf{Autonomía de 7+ horas} con configuración de baterías 18650, suficiente para jornadas laborales completas.
  \item \textbf{Latencia de 2.8 segundos} compatible con la velocidad de respuesta operativa requerida.
  \item \textbf{Instalación no invasiva} que no interfiere con sistemas del fabricante.
  \item \textbf{Protocolo robusto} con buffer automático ante desconexiones.
  \item \textbf{Precisión diferencial de $\pm$2\%} mediante validación cruzada ida-retorno.
  \item \textbf{Evidencia digital persistente} incluso si el ladrón reconecta mangueras (historial de eventos en memoria).
\end{itemize}

Como proyección futura, se identifica la necesidad de: (i)~expandir las pruebas de campo a flotas múltiples para validar el modelo de aprendizaje automático (RPM--consumo esperado), (ii)~mejorar el sellado industrial de los nodos a IP67 para resistir la corrosión prolongada por diésel, (iii)~optimizar el consumo energético del módulo 4G en estado de reposo mediante configuración de power management, y (iv)~implementar certificaciones de seguridad industrial para comercialización. Adicionalmente, se propone investigar la integración de machine learning para optimización automática de umbrales ($\delta$ y $\epsilon$) basada en patrones de consumo específicos por modelo de vehículo.

%% ─────────────────────────────────────────────
\section*{AGRADECIMIENTOS}
%% ─────────────────────────────────────────────

Los autores expresan su agradecimiento al Instituto Tecnológico Superior del Occidente del Estado de Hidalgo por las facilidades otorgadas para el uso de los laboratorios de Internet de las Cosas. Asimismo, se reconoce profundamente a la profesora Eunice Santiago Manzano por su orientación metodológica en la estructuración de este artículo; al profesor Pedro Jhoan Salazar Perez por su asesoría técnica especializada en la selección e instrumentación de sensores; y al profesor Saul Isai Soto Ortis por su dirección, validación y aprobación técnica durante el diseño del prototipo.

%% ─────────────────────────────────────────────
\section*{REFERENCIAS}
%% ─────────────────────────────────────────────
\small
\begin{enumerate}

\item Atzori, L., Iera, A., \& Morabito, G. (2010). The Internet of Things: A survey. \textit{Computer Networks}, 54(15), 2787--2805.

\item Al-Fuqaha, A., Guizani, M., Mohammadi, M., Aledhari, M., \& Ayyash, M. (2015). Internet of Things: A survey on enabling technologies. \textit{IEEE Communications Surveys \& Tutorials}, 17(4), 2347--2376.

\item Zanella, A., Bui, N., Castellani, A., Vangelista, L., \& Zorzi, M. (2014). Internet of Things for smart cities. \textit{IEEE Internet of Things Journal}, 1(1), 22--32.

\item Gubbi, J., Buyya, R., Marusic, S., \& Palaniswami, M. (2013). Internet of Things: A vision, architectural elements, and future directions. \textit{Future Generation Computer Systems}, 29(7), 1645--1660.

\item Vermesan, O., \& Friess, P. (2014). \textit{Internet of Things: From research and innovation to market deployment}. River Publishers.

\item Da Xu, L., He, W., \& Li, S. (2014). Internet of Things in industries: A survey. \textit{IEEE Transactions on Industrial Informatics}, 10(4), 2233--2243.

\item Ray, P. (2018). A survey on Internet of Things architectures. \textit{Journal of King Saud University – Computer and Information Sciences}, 30(3), 291--319.

\item Borgia, E. (2014). The Internet of Things vision: Key features, applications and open issues. \textit{Computer Communications}, 54, 1--31.

\item Raza, U., Kulkarni, P., \& Sooriyabandara, M. (2017). Low power wide area networks: An overview. \textit{IEEE Communications Surveys \& Tutorials}, 19(2), 855--873.

\item Centenaro, M., Vangelista, L., Zanella, A., \& Zorzi, M. (2016). Long-range communications in unlicensed bands. \textit{IEEE Wireless Communications}, 23(5), 60--67.

\item Palattella, M. R., et al. (2016). Internet of Things in the 5G era. \textit{IEEE Communications Magazine}, 54(12), 140--147.

\item Li, S., Xu, L., \& Zhao, S. (2018). The Internet of Things: A survey. \textit{Information Systems Frontiers}, 17(2), 243--259.

\item Lin, J., Yu, W., Zhang, N., Yang, X., Zhang, H., \& Zhao, W. (2017). A survey on Internet of Things. \textit{IEEE Internet of Things Journal}, 4(5), 1125--1142.

\item Perera, C., Zaslavsky, A., Christen, P., \& Georgakopoulos, D. (2014). Context aware computing for the Internet of Things. \textit{IEEE Communications Surveys \& Tutorials}, 16(1), 414--454.

\item Madakam, S., Ramaswamy, R., \& Tripathi, S. (2015). Internet of Things: A literature review. \textit{Journal of Computer and Communications}, 3(5), 164--173.
\item u-blox. (2013). NEO-6 and NEO-6M receiver documentation. \textit{u-blox AG}, SWISS.

\item Semtech. (2015). LoRa and LoRaWAN: A technical overview. \textit{Semtech Corporation}.

\item National Instruments. (2019). Wireless sensor networks in IoT applications. \textit{NI Technical Publications}.

\item IEEE Standards Association. (2014). IEEE 802.11 Wireless LAN standards for IoT. \textit{IEEE Standard}, 802.11.

\item Cherry Semiconductor. (2018). Reed Switch Reliability in IoT Applications. \textit{Cherry Electronic Components}, Datasheet Series.

\item Panasonic Industrial. (2019). Ultrasonic flow sensors: Principles and applications in fuel monitoring systems. \textit{Panasonic Technical Publications}.
\end{enumerate}

\end{document}